\documentclass[11pt]{article}

\usepackage{amsmath,amssymb}
\usepackage{xurl}
\usepackage[colorlinks=true,linkcolor=blue,citecolor=blue,urlcolor=blue]{hyperref}

\input{command}

\title{The Schmidt-rank set in the spectral $n$-positivity criterion
  (Wolf Prop.~3.2)}
\author{}
\date{2026-08-13}

\begin{document}
\maketitle
\gapnote{scope-restriction}{resolved}

\section{The assertion}

Let $\tau\in\mathcal M_{D'}\otimes\mathcal M_D$ be Hermitian, with spectral
decomposition
\begin{align}
    \tau&=\sum_i\nu_i\ketbra{\phi_i}{\phi_i},
    \label{eq:wolf_prop_3_2_schmidt_rank_reading_spectral_decomposition}\\
    \rho_i&=\tr_2\!\left(\ketbra{\phi_i}{\phi_i}\right).
    \label{eq:wolf_prop_3_2_schmidt_rank_reading_reduced_density}
\end{align}
Here the eigenvectors $\phi_i$ are normalized, and $\rho_i$ is their reduced
density matrix on the first tensor factor. For a positive semidefinite matrix
$\rho$, let $\lVert\rho\rVert_{(n)}$ denote its Ky--Fan $n$-norm, the sum of
its $n$ largest eigenvalues.

Let $\nu_0$ be the smallest positive eigenvalue of $\tau$, and let $\nu$ be
its largest eigenvalue. Proposition~3.2 of Ref.~\cite{Wolf2012Quantum} states
the lower bound
\begin{align}
    \inf_{\psi}\bra{\psi}\tau\ket{\psi}
        &\geq \nu_0+
        \sum_{i:\nu_i\leq0}(\nu_i-\nu_0)\lVert\rho_i\rVert_{(n)}.
    \label{eq:wolf_prop_3_2_schmidt_rank_reading_source_lower_bound}
\end{align}
This is Equation~(3.7) of Ref.~\cite{Wolf2012Quantum}. If exactly one
eigenvalue $\nu_-$ is non-positive and $\rho_-$ is the reduced density matrix
of a corresponding normalized eigenvector, the proposition also states
\begin{align}
    \inf_{\psi}\bra{\psi}\tau\ket{\psi}
        &\leq \nu+(\nu_--\nu)\lVert\rho_-\rVert_{(n)}.
    \label{eq:wolf_prop_3_2_schmidt_rank_reading_source_upper_bound}
\end{align}
This is Equation~(3.8) of Ref.~\cite{Wolf2012Quantum}. The local source is
\path{Notes/WolfNoteTexSource/ch03_positive_not_completely.tex}, line~153,
including Equations~(3.7) and~(3.8). In both statements, the source says that
the infimum is over normalized vectors ``of Schmidt rank $n$.''

\section{The point at issue}

For $1\leq n\leq\min(D',D)$, define
\begin{align}
    S_{\leq n}
        &=\left\{\psi:\lVert\psi\rVert=1,\
            \operatorname{SR}(\psi)\leq n\right\},
    \label{eq:wolf_prop_3_2_schmidt_rank_reading_bounded_rank_set}\\
    S_{=n}
        &=\left\{\psi:\lVert\psi\rVert=1,\
            \operatorname{SR}(\psi)=n\right\}.
    \label{eq:wolf_prop_3_2_schmidt_rank_reading_exact_rank_set}
\end{align}
Then
\begin{align}
    S_{=n}&\subseteq S_{\leq n}.
    \label{eq:wolf_prop_3_2_schmidt_rank_reading_set_inclusion}
\end{align}
The phrase ``Schmidt rank $n$'' can denote either set.

The proof of Proposition~3.1 in Ref.~\cite{Wolf2012Quantum} fixes the intended
reading. It expresses $n$-positivity of $T$ by requiring
\begin{align}
    (T\otimes\operatorname{id}_n)
        \!\left(\ketbra{\phi}{\phi}\right)&\geq0
    \label{eq:wolf_prop_3_2_schmidt_rank_reading_n_positive_test}
\end{align}
for every $\phi\in\C^D\otimes\C^n$. Under the embedding
$\C^n\subseteq\C^D$, these vectors have Schmidt rank at most $n$, not
necessarily exactly $n$. Vectors of smaller Schmidt rank remain in the
quantifier. The source makes the same identification later when it writes
\begin{align}
    \ket{\psi}
        &=(\Id\otimes X)^\dagger\ket{\widetilde\psi},
    \label{eq:wolf_prop_3_2_schmidt_rank_reading_factorization}\\
    \operatorname{rank}(X)&=n,
    \label{eq:wolf_prop_3_2_schmidt_rank_reading_factor_rank}
\end{align}
and calls $\psi$ a vector of Schmidt rank $n$, although this construction only
implies $\operatorname{SR}(\psi)\leq n$. The relevant local source is
\path{Notes/WolfNoteTexSource/ch03_positive_not_completely.tex},
Proposition~3.1 and its proof.

This bounded-rank interpretation also carries the operational content:
$n$-positivity of a Hermitian map is equivalent to nonnegativity of the
corresponding Choi quadratic form on vectors in $S_{\leq n}$.

\section{Comparison of the two readings}

The lower bound is pointwise. Every $\psi\in S_{\leq n}$ satisfies
\begin{align}
    \bra{\psi}\tau\ket{\psi}
        &\geq \nu_0+
        \sum_{i:\nu_i\leq0}
        (\nu_i-\nu_0)\lVert\rho_i\rVert_{(n)},
    \label{eq:wolf_prop_3_2_schmidt_rank_reading_pointwise_lower_bound}
\end{align}
because
\begin{align}
    \left|\braket{\phi_i}{\psi}\right|^2
        &\leq\lVert\rho_i\rVert_{(n)}.
    \label{eq:wolf_prop_3_2_schmidt_rank_reading_overlap_bound}
\end{align}
By (\ref{eq:wolf_prop_3_2_schmidt_rank_reading_set_inclusion}), the same
pointwise estimate holds on $S_{=n}$. The bounded-rank version of the lower
bound is therefore stronger.

The upper bound uses a vector that maximizes the overlap with $\phi_-$. The
standard construction from the $n$ largest eigenvalues of $\rho_-$ has
Schmidt rank
\begin{align}
    \operatorname{SR}(\psi_{\max})
        &=\min\!\left(n,\operatorname{rank}(\rho_-)\right),
    \label{eq:wolf_prop_3_2_schmidt_rank_reading_witness_rank}\\
    \left|\braket{\phi_-}{\psi_{\max}}\right|^2
        &=\lVert\rho_-\rVert_{(n)}.
    \label{eq:wolf_prop_3_2_schmidt_rank_reading_maximal_overlap}
\end{align}
If $\operatorname{rank}(\rho_-)<n$, then $\psi_{\max}\in S_{\leq n}$ but
$\psi_{\max}\notin S_{=n}$. In this case,
\begin{align}
    \nu+(\nu_--\nu)\lVert\rho_-\rVert_{(n)}
        &=\nu_-.
    \label{eq:wolf_prop_3_2_schmidt_rank_reading_upper_value}
\end{align}
No vector in $S_{=n}$ attains this value: an attaining vector would belong to
the $\nu_-$ eigenspace, which is spanned by $\phi_-$ and contains no vector of
Schmidt rank $n$.

Nevertheless, $\phi_-$ is a limit of vectors in $S_{=n}$ whenever
$n\leq\min(D',D)$. Continuity of the quadratic form therefore gives
\begin{align}
    \inf_{\psi\in S_{\leq n}}\bra{\psi}\tau\ket{\psi}
        &=\inf_{\psi\in S_{=n}}\bra{\psi}\tau\ket{\psi}.
    \label{eq:wolf_prop_3_2_schmidt_rank_reading_equal_infima}
\end{align}
Thus the two readings produce the same infimum bounds, although the exact-rank
upper bound may be approached only by a limiting sequence.

\section{The formalized statement}

The formalization uses the normalized vectors of Schmidt rank at most $n$ and
states both inequalities for the infimum of their expectation values.
\footnote{Formal declarations:
    \leanid{Matrix.schmidtRankLEExpectations},
    \leanid{Matrix.IsHermitian.le_sInf_schmidtRankLEExpectations} for
    Equation~(3.7) of Ref.~\cite{Wolf2012Quantum},
    \leanid{Matrix.IsHermitian.sInf_schmidtRankLEExpectations_le} for
    Equation~(3.8) of Ref.~\cite{Wolf2012Quantum}, and
    \leanid{Matrix.IsHermitian.sInf_schmidtRankLEExpectations_top} for the top
    index, all in
    \doclink{QICLean/Channel/NPositivitySpectralCriterion}.}
No hypothesis absent from Ref.~\cite{Wolf2012Quantum} is used.

The identity
(\ref{eq:wolf_prop_3_2_schmidt_rank_reading_equal_infima}) is not formalized.
A proof requires normalized vectors of exact Schmidt rank $n$ and a
rank-raising perturbation. For a normalized vector $\psi$ and a product vector
$u\otimes v$ orthogonal to its Schmidt vectors, set
\begin{align}
    \psi_\varepsilon
        &=\sqrt{1-\varepsilon^2} \psi+\varepsilon\,u\otimes v.
    \label{eq:wolf_prop_3_2_schmidt_rank_reading_rank_raising_perturbation}
\end{align}
The expectation changes according to
\begin{align}
    \bra{\psi_\varepsilon}\tau\ket{\psi_\varepsilon}
        -\bra{\psi}\tau\ket{\psi}
        &=
        2\varepsilon\sqrt{1-\varepsilon^2}
        \Re\!\left(\bra{\psi}\tau\ket{u\otimes v}\right)
    \label{eq:wolf_prop_3_2_schmidt_rank_reading_expectation_cross_term}\\
        &\quad+
        \varepsilon^2
        \left(\bra{u\otimes v}\tau\ket{u\otimes v}
            -\bra{\psi}\tau\ket{\psi}\right).
    \label{eq:wolf_prop_3_2_schmidt_rank_reading_expectation_quadratic_term}
\end{align}
The change is $O(\varepsilon)$ in general. The cross term in
(\ref{eq:wolf_prop_3_2_schmidt_rank_reading_expectation_cross_term}) vanishes
only under an additional orthogonality condition involving $\tau$, which the
argument does not require. Iterating such perturbations and taking
$\varepsilon\to0$ proves the nontrivial direction of
(\ref{eq:wolf_prop_3_2_schmidt_rank_reading_equal_infima}).

\section{Verdict}

The source phrase ``vectors of Schmidt rank $n$'' means vectors of Schmidt rank
at most $n$, as the proof of Proposition~3.1 in
Ref.~\cite{Wolf2012Quantum} shows. The formalized lower and upper bounds are
therefore faithful to Equations~(3.7) and~(3.8) of
Ref.~\cite{Wolf2012Quantum}, and no additional hypothesis is introduced.

The residual scope restriction concerns only the exact-rank reading of the
upper bound. The lower bound follows immediately from the bounded-rank result,
whereas the upper bound requires the limiting argument leading to
(\ref{eq:wolf_prop_3_2_schmidt_rank_reading_equal_infima}) rather than an
exact-rank witness. That limiting argument is not formalized and is not needed
for the source-facing bounded-rank statement. The Lean statement corresponding
to Equation~(3.8) of Ref.~\cite{Wolf2012Quantum} carries a
``Scope restriction (Schmidt rank at most $n$)'' marker pointing to this note.

\bibliographystyle{alpha}
\bibliography{references}

\end{document}
